\section*{第 \ref{ch:b2:vertebrate-development} 章——脊椎动物的胚胎发育}

\begin{solution}{exo:b2:vertebrate-development:1}
\omterm{def:b2:vertebrate-development:egg}{卵裂}：快速的有丝分裂，不伴随生长，产生\omterm{def:b2:vertebrate-development:blastula}{囊胚}。
\omterm{def:b2:vertebrate-development:blastula}{囊胚}：由小细胞围成的空心球或圆盘，围绕
\omterm{def:b2:vertebrate-development:blastula}{囊胚腔}。原肠胚形成：使中胚层和
内胚层进入内部的细胞运动，产生带有肠腔
和脊索的三胚层原肠胚。\omterm{prop:b2:vertebrate-development:neurulation}{神经胚形成}：背侧
外胚层折叠成\omterm{prop:b2:vertebrate-development:neurulation}{神经管}，同时\omterm{prop:b2:vertebrate-development:neurulation}{体节}分节，
产生具有脊椎动物身体构型的神经胚。
\end{solution}

\begin{solution}{exo:b2:vertebrate-development:2}
外胚层：表皮、脑和脊髓、外周神经
（\omterm{prop:b2:vertebrate-development:neurulation}{神经嵴}）、晶状体和色素细胞。中胚层：脊索、肌肉、
骨和软骨、肾脏、心脏和血液、真皮。内胚层：
肠道内衬、肝脏、胰腺、肺、甲状腺。
\end{solution}

\begin{solution}{exo:b2:vertebrate-development:3}
\omterm{def:b2:vertebrate-development:egg}{卵黄}少：完全、近似均等的\omterm{def:b2:vertebrate-development:egg}{卵裂}，原肠胚形成通过
空心球中圆形\omterm{prop:b2:vertebrate-development:gastrulation}{胚孔}的内卷进行（蛙：
完全但不均等，富含\omterm{def:b2:vertebrate-development:egg}{卵黄}的\omterm{def:b2:vertebrate-development:egg}{植物极}细胞大而分裂慢）。\omterm{def:b2:vertebrate-development:egg}{卵黄}多
（鸡）：\omterm{def:b2:vertebrate-development:egg}{卵裂}局限于\omterm{def:b2:vertebrate-development:egg}{卵黄}顶部的圆盘中，
原肠胚形成通过平盘中的裂缝（原条）进行，
各胚层在\omterm{def:b2:vertebrate-development:egg}{卵黄}上铺展，而不是包围
一个腔。
\end{solution}

\begin{solution}{exo:b2:vertebrate-development:4}
诱导是指信号使有\omterm{def:b2:vertebrate-development:induction}{感受态}的组织
命运发生改变，信号来自邻近组织。\omterm{def:b2:vertebrate-development:egg}{植物极}细胞在上方
的边缘区诱导中胚层；\omterm{def:b2:vertebrate-development:induction}{组织者}（背唇中胚层）诱导
背侧外胚层中的神经板；视泡诱导
头部外胚层形成晶状体。
\end{solution}

\begin{solution}{exo:b2:vertebrate-development:5}
$2^{12} = 4096$ 个细胞，半径为 $1/2^{4} = \qty{62.5}{\micro m}$；
表面积变为原来的 $2^{4} = 16$ 倍。额外的表面积是
上皮的膜，原肠胚形成时将把它折叠成皮肤和
肠道，也是\omterm{def:b2:vertebrate-development:blastula}{囊胚}与
外界交换的面积。
\end{solution}

\begin{solution}{exo:b2:vertebrate-development:6}
$10^{-4}\times 2^{k} = 0.4$：$2^{k} = 4000$，$k = 12$。
若 DNA 量为四倍，比值起始为 $4\times 10^{-4}$，达到 $0.4$
时对应 $2^{k} = 1000$：第十次分裂，比原来早两次。
\end{solution}

\begin{solution}{exo:b2:vertebrate-development:7}
$k = \ln 2/600 = \qty{1.16e-3}{s^{-1}}$；$\lambda = \sqrt{2\times
10^{-12}/1.16\times 10^{-3}} = \qty{42}{\micro m}$。阈值位于
$x = \lambda\ln 2 = \qty{29}{\micro m}$ 和 $\lambda\ln 10 =
\qty{96}{\micro m}$ 处，条带分别为 29、67 和其余部分。
\end{solution}

\begin{solution}{exo:b2:vertebrate-development:8}
源加倍使两个边界都向外移动 $\lambda\ln
2 = \qty{29}{\micro m}$：第一条条带加倍，第二条
条带的宽度保持不变。$k$ 加倍使 $\lambda$ 除以 $\sqrt 2$
（\qty{29}{\micro m}）；由于 $c_0 = j/\sqrt{Dk}$，$c_0$ 也被
$\sqrt 2$ 除，于是边界移至 $\lambda'\ln(c_0'/T) =
29\times(0.69 - 0.35) = \qty{10}{\micro m}$ 和 $29\times(2.30 -
0.35) = \qty{57}{\micro m}$——两条条带均缩小。
\end{solution}

\begin{solution}{exo:b2:vertebrate-development:9}
将一个无色素的蝾螈原肠胚的背唇移植到
有色素的原肠胚腹侧；背唇卷入其中，宿主形成了
第二个胚胎——\omterm{prop:b2:vertebrate-development:neurulation}{神经管}、\omterm{prop:b2:vertebrate-development:neurulation}{体节}、肠道——与第一个相连。
色素差异表明，第二条体轴由宿主细胞构成，
因而移植体诱导了宿主细胞，而非自身构建
体轴。这排除了移植体自分化这一解释，
并表明一小块区域携带着
组织整个体轴的信号。
\end{solution}

\begin{solution}{exo:b2:vertebrate-development:10}
(a) 腹侧皮肤：早期外胚层有\omterm{def:b2:vertebrate-development:induction}{感受态}，但尚未被诱导，
因而随新的周围环境而变化。(b) 神经组织：至原肠胚后期，
它已被诱导，且已决定命运。(c) 神经板：
原肠胚后期的腹侧外胚层仍有\omterm{def:b2:vertebrate-development:induction}{感受态}，脊索
诱导它。(d) 无背侧结构：一团腹侧组织，
因为没有东西诱导中胚层走向背侧命运，
或诱导外胚层走向神经命运。
\end{solution}

\begin{solution}{exo:b2:vertebrate-development:11}
在 \qty{18}{\celsius} 下：$90 + 11\times 30 = \qty{420}{min}$，7 小时。
在 \qty{10}{\celsius} 下：\qty{17.5}{h}。该转换发生在第
两种情况下都在第十二次分裂时发生，因为触发它的是
DNA 与细胞质的比值，而该比值取决于分裂次数，
与分裂持续多长时间无关：计时的是倍增次数，而非
小时。
\end{solution}

\begin{solution}{exo:b2:vertebrate-development:12}
在调整型发育中，细胞通过来自邻近细胞的诱导获得
其命运，因此胚胎在损失后可以重建（由分裂的卵形成双胚，
或移植区域的再生）；在镶嵌型发育中，
命运随定位的细胞质决定因子遗传，
因此每个\omterm{def:b2:vertebrate-development:egg}{卵裂}球只形成自身的部分，分裂的
胚胎形成两个半幼体。二者的区别只在于程度：
蛙的灰色新月是一种决定因子，而镶嵌型胚胎
后期才使用诱导；每个胚胎都兼有遗传与“对话”，
脊椎动物更长时间地依赖“对话”。
\end{solution}

\begin{solution}{pb:b2:vertebrate-development:1}
\textbf{1.} 卵 $\tfrac43\pi(0.7)^{3} = \qty{1.44}{mm^3}$；细胞核
$\tfrac43\pi(0.03)^{3} = \qty{1.1e-4}{mm^3}$；比值 $8\times
10^{-5}$。
\textbf{2.} $2^{k} = 4000$：$k = 12$，4096 个细胞。
\textbf{3.} $75 + 11\times 30 = \qty{405}{min}$，约 6.75 小时。
\textbf{4.} 半径 $0.7/2^{4} = \qty{44}{\micro m}$；比值 $8\times
10^{-5}\times 4096 = 0.32$——与普通细胞相同：细胞核
与细胞质之比已追平。
\textbf{5.} $2^{4} = 16$。
\textbf{6.} 约 $4095\,c$ 的 DNA；卵中储存的
组蛋白、聚合酶和核苷酸足以支持十二次分裂，
因此不需要转录。
\textbf{7.} 合子基因开启，分裂减慢并失去
同步性，细胞变得可运动。注入额外的 DNA 会使
转换提前相当于过量 DNA 所对应的分裂次数，
这是按分裂次数或时间计数无法解释的。
\textbf{8.} 与精子进入点相对，位于灰色新月处，
由受精后卵皮层的转动所固定。
\textbf{9.} 背侧中胚层先（脊索前板，随后为脊索）
沿中线迁至原肠顶部；侧中胚层
（\omterm{prop:b2:vertebrate-development:neurulation}{体节}、侧板）位于其旁；内胚层衬于
原肠内壁；外胚层经外包作用覆盖外侧。
\textbf{10.} $\qty{2000}{\micro m}/\qty{600}{min} = \qty{3.3}{\micro
m/min}$：为成纤维细胞的三倍——协同推进的细胞层
比单个细胞移动得更快。
\textbf{11.} 细胞运动需要合子基因产物（新的
粘附分子、运动能力）以及放慢的细胞周期；在
转换之前，细胞只能分裂。
\textbf{12.} 所有脊索动物在某一阶段都有脊索，而
脊椎动物围绕它构建脊柱；\omterm{prop:b2:vertebrate-development:neurulation}{神经管}
由它诱导产生，是身体构型的结果而非
起源。
\textbf{13.} 若无瓶状细胞收缩，就不会形成原口，
也无内卷，胚胎仍是一个各层位置未变的球，
没有原肠、脊索，也没有轴。
\textbf{14.} $k = \ln 2/1200 = \qty{5.8e-4}{s^{-1}}$；$\lambda =
\sqrt{10^{-12}/5.8\times 10^{-4}} = \qty{42}{\micro m}$。
\textbf{15.} $x_1 = 42\ln 2.5 = \qty{38}{\micro m}$；$x_2 = 42\ln 12.5
= \qty{105}{\micro m}$。
\textbf{16.} 条带宽度为 \qty{38}{\micro m}（2.5 个细胞）、\qty{67}{\micro
m}（4.5 个细胞）以及其余部分。
\textbf{17.} $x_1 = 42\ln 1.25 = \qty{9}{\micro m}$；$x_2 = 42\ln 6.25
= \qty{77}{\micro m}$：两条边界都向内移动 $\lambda\ln 2 =
\qty{29}{\micro m}$；第一条带由 38 缩小到 9，中间
条带宽度不变，外侧条带变宽。
\textbf{18.} 读取之前移动，则变为命运 A——它读取
所在位置的浓度；读取之后移动，则保持命运 B——它
已被决定。
\textbf{19.} $x_T = \lambda\ln(c_0/T)$：源的变化为
$f$ 倍时，每条边界移动 $\lambda\ln f$，因此两倍误差
只使图案移动 $0.7\lambda$，仅为跨越数个 $\lambda$ 的场的
一小部分——图案对\omterm{thm:b2:vertebrate-development:gradient}{形态发生素}的
确切用量几乎不敏感。
\textbf{20.} 背唇移植到腹侧：形成第二条轴。腹侧边缘
区移植到背侧：它被周围环境重新定型，形成
背侧组织；没有额外结构。
\textbf{21.} 在神经胚期，外胚层已不再具有\omterm{def:b2:vertebrate-development:induction}{感受态}：
移植物形成脊索，却不诱导第二个\omterm{prop:b2:vertebrate-development:neurulation}{神经管}。
\textbf{22.} 沿新月平面切开：得到两个完整胚胎；
垂直于该平面切开：一个胚胎及一团腹侧组织。
\textbf{23.} 单独时：形成表皮球。与\omterm{def:b2:vertebrate-development:egg}{植物极}细胞共同时：中胚层
——肌肉、脊索——在动物帽中被诱导产生。
\textbf{24.} 为证明第二条轴是在宿主细胞中被诱导的，
而非由移植物自身构建；若移植物与宿主不可区分，
则无法排除自分化的可能。来自不同物种的移植同样成功（他们用了两种
蝾螈），这正是使细胞可以区分的
原因。
\textbf{25.} 第十二次\omterm{def:b2:vertebrate-development:egg}{卵裂}，4096 个细胞，约 6.75 小时；条带
宽为 38 与 \qty{67}{\micro m}，以及超过 \qty{105}{\micro
m} 的全部区域。
\end{solution}
